% Springer LNCS manuscript, main text (references included). Appendices are in supplementary.tex.
% Code and per-sample results: https://github.com/tsan2711/DACNTT_2026
\documentclass[runningheads]{llncs}

\usepackage[T1]{fontenc}
\usepackage{graphicx}
\usepackage{booktabs}
\usepackage{amsmath}
\usepackage{array}
\usepackage{xcolor}
\usepackage{hyperref}
\usepackage{xr-hyper}
\externaldocument{supplementary}
\renewcommand\UrlFont{\color{blue}\rmfamily}
\urlstyle{rm}

\begin{document}

\title{Mathematics or the Verifier? A Null Result for One
Verifier-Filtered Self-Training Recipe on a 1.5B Model}

\titlerunning{Mathematics or the Verifier?}

\author{Tan Sang Nguyen\inst{1} \and Duy Anh Tran\inst{1}}
\authorrunning{T. S. Nguyen and D. A. Tran}
\institute{Ton Duc Thang University, Ho Chi Minh City, Vietnam\\
\email{52300057@student.tdtu.edu.vn, 52300006@student.tdtu.edu.vn}}

\maketitle

\begin{abstract}
Self-training and reinforcement learning with verifiable rewards report large
gains for small language models, but a rising score does not show that a
model learned mathematics rather than what its verifier accepts. We study the
cheapest form of the loop, generate--verify--select--train with offline LoRA
SFT, on Qwen2.5-1.5B-Instruct with the rule-based verifier
\texttt{math-verify}, on held-out MATH-500 and GSM8K exams (150 problems; $500$
in a final round) at four learning rates on MATH-500 and two on GSM8K. No
pass@1 or pass@$k$ gain has a problem-level bootstrap interval excluding zero;
in the final round, with a revised (not validated) recipe, pass@1 falls $3.2$--$4.8$ points
in each of three seeds (mean $-4.1$; two of three intervals exclude zero). The design detects only gains of about $10$ points of
pass@1 ($6$ in the final round; $80\%$ power, intervals without seed
variance, so likely more). A positive gate fails: training on reference solutions
lowers pass@1 by $32.2$ points, which we did not diagnose (a pipeline defect would also produce it), so the null
concerns this recipe only and cannot show that self-training fails. Accuracies are graded by the
verifier that selects the data (the final round also by a string-match grader), so neither the mathematics nor the verifier question is answered. As a post hoc audit we track the share of the loop's
own correct solutions that a storage defect in our pipeline rejects (about one
in six). It falls by $1.5$--$2.6$ points with intervals reaching zero; this fall is tied to the
falling \verb|\boxed| share by an identity, so it is no evidence of adaptation. Two injected-bias runs, one whose rule is visible in the gold answer, show no selective avoidance but are not
valid positive controls (one cannot elicit avoidance by construction, the other is a single run without $r_t$), so the audit is uninformative about adaptation.
Two controls, patching a rarer defect
and removing the filter, show no consistent difference (no filter is lower on pass@$k$, at $3$--$4\times$ the dose) but are dose-confounded
or too small to detect one. Scores before the final round use the verifier with the defect.

\keywords{self-training \and self-improvement \and verifier bias \and
mathematical reasoning \and small language models \and verifiable rewards
\and null result.}
\end{abstract}

\section{Introduction}
\label{sec:intro}

A common recipe for improving the mathematical reasoning of a small language
model (SLM) \emph{without human annotation} is to let the model generate its own
solutions, keep those an automated verifier judges correct, and fine-tune on
them, for several rounds. We call this the \emph{Generate--Verify--Train} (GVT)
loop; it underlies STaR~\cite{ref_star}, ReST~\cite{ref_rest},
ReST-EM~\cite{ref_restem} and rejection-sampling fine-tuning~\cite{ref_rft}.
But the verifier never reads the reasoning: a rule-based verifier such as
\texttt{math-verify}\footnote{\url{https://github.com/huggingface/math-verify},
version 0.9.0 with antlr4 4.13.2 in this work.} only extracts a final answer and
compares it with the gold, and it rejects some correct answers written in forms
it does not parse (Section~\ref{sec:verifier}). The training signal is ``was the
final string accepted'', not ``were the steps right''.

The same gap runs through reinforcement learning with verifiable rewards
(RLVR), which reports large gains at this scale~\cite{ref_oneshot,ref_openrs,ref_simplerl}, though only $8.6$ of one study's
$37.6$ points lie beyond format correction. Shao et al.~\cite{ref_spurious} find that
GRPO with \emph{random} rewards raises Qwen2.5-Math-7B by $21.4$ points on
MATH-500, against $29.1$ with ground-truth rewards. A rising score therefore
does not by itself show that a model has learned to solve more problems,
rather than to surface behaviour it already had or to produce what its reward
accepts. In RLVR, per-step false-negative rates on the policy's own rollouts
have been tracked~\cite{ref_delayplateau,ref_tinyv}; we study the question in a cheap and common
form of the loop, offline rejection-sampling SFT, where a false reject
removes a training example (and, if bound to a problem, the problem), and apply a per-round false-reject
measurement there, post hoc, for one defect of our pipeline.

We ask: \emph{after several GVT rounds, is the model better at mathematics, or
has it only learned to write what the verifier accepts?} We fixed four measures
in advance: (i)~pass@1; (ii)~pass@$k$; (iii)~output style, chiefly the share of
generations containing \verb|\boxed|; and (iv)~per-problem transitions
(problems \emph{added} and \emph{lost} between rounds). The pre-set reading rule
was that ``learning the verifier's style'' requires pass@1 to rise, pass@$k$ to
stay flat and style to converge at once. In hindsight the rule is one-sided: it
can only fire when accuracy rises (Section~\ref{sec:answer}). The higher learning rates, the
controls, the round count and $k$ at those rates, sample accuracy as primary
measure, the termination measures and the re-scoring for a second verifier
defect, and the final 500-problem round
(Section~\ref{sec:final}) were all added after seeing earlier runs; all are post hoc.

We detect no improvement at any learning rate, though gains below about $10$
points of pass@1 or $3.5$--$5.5$ of sample accuracy could have been missed
(intervals omit training-seed variance), and no arm of this recipe improves, so the result concerns
the recipe. The second branch is untested: no run varies the frequent verifier defect, scores are
graded by the selecting verifier (one string-match check on the final round), and the post hoc audit of
Section~\ref{sec:audit} is uninformative about adaptation. What appears,
most clearly at the highest rate, is lower accuracy on the solutions the model
does finish, though on one seed the cap share carries more of it (Section~\ref{sec:anatomy}). A sharp pooled decline is absent in the isolated runs at $10^{-5}$, which train on about $4.5\times$ fewer examples
(Section~\ref{sec:pooled}); dose and pooling are confounded.

\paragraph{Contributions.}
\begin{itemize}
\item For a 1.5B model, LoRA SFT and one rule-based verifier, no gain is
detectable on a held-out exam at any learning rate we tried (a coarse grid that also changes $k$ and rounds), with
bootstrap intervals and the smallest effects the design could detect. No arm
improves, not even training on reference solutions, so this is a null for the
recipe, not a statement about self-training.
\item A post hoc per-round false-reject audit of one verifier defect (the RLVR
tracking~\cite{ref_delayplateau,ref_tinyv} applied to offline SFT; not a novel method). The rejected
share falls by $1.5$--$2.6$ points with intervals reaching zero, tied to the falling \verb|\boxed| share by
an identity; the audit is uninformative about adaptation (Section~\ref{sec:audit}).
\item A decomposition of an accuracy decline under a fixed token cap into a
finished-accuracy part and a cap-share part: at the highest rate, finished
generations lose $5.5$--$10.6$ points against the untrained model on the same
problem mix (A, B and no-filter runs; train-on-gold loses about $26$), and the cap-hit share rises only at that rate.
\end{itemize}

\section{Related Work}
\label{sec:related}

\paragraph{Self-training and its limits.}
STaR~\cite{ref_star} introduced the generate--filter--retrain loop, and
ReST~\cite{ref_rest}, ReST-EM~\cite{ref_restem} and RFT~\cite{ref_rft} formalise
it; ReST-EM, which like us restarts from the base model each iteration, reports
diminishing gains after the first round. Yue et al.~\cite{ref_rllimit} show that RL
raises pass@1 while the base model wins at large pass@$k$, and Teacher-Free
Self-Training~\cite{ref_teacherfree} concludes that self-training ``amplifies but
does not compound''. Training repeatedly on model-generated data erodes
distribution tails~\cite{ref_collapse}, and a narrow selector can accelerate
collapse~\cite{ref_selectionbias}.

\paragraph{RL with verifiable rewards at small scale.}
RLVR with GRPO reports large gains near our model size: one training example
lifts Qwen2.5-Math-1.5B on MATH-500 from $36.0\%$ to $73.6\%$~\cite{ref_oneshot},
$7{,}000$ examples lift DeepSeek-R1-Distill-Qwen-1.5B on AMC23 from $63\%$
to $80\%$~\cite{ref_openrs}, and zero-RL works on open base models from $0.5$B~\cite{ref_simplerl}. These studies improve
accuracy, which our SFT loop does not; of Wang et al.'s $37.6$ points only $8.6$ lie beyond format correction. Shao et al.~\cite{ref_spurious} show that it need not be
reward learning (random rewards also help, via amplified pre-existing behaviour). We do not compete with RLVR on accuracy; the audit of
Section~\ref{sec:audit} asks whether a loop has adapted to one
verifier defect, and we apply it only to SFT.

\paragraph{Verifier error.}
TinyV~\cite{ref_tinyv} finds that rule-based verifiers reject many correct
solutions; From Accuracy to Robustness~\cite{ref_robustness} reports an average
recall of $86\%$, and $0.95$ for \texttt{math-verify} on their MATH subset.
Imperfect Verifiers~\cite{ref_imperfect} models fixed false-accept and
false-reject rates. Egashira et al.~\cite{ref_delayplateau} find in RLVR that
systematic false negatives behave like random noise whereas false positives can
cause plateau or collapse. For format-based false negatives they see the policy
first unlearn the triggering behaviour. Our setting differs in
paradigm (offline SFT with a fresh adapter per round) and task, and our
controls cannot test their findings (Section~\ref{sec:control}).

\paragraph{Auditing self-improvement.}
Xu et al.~\cite{ref_phantom} audit LoRA self-training of Qwen3-8B against a
frozen control and identify seven measurement failures, including training-seed
variance and a fixed token cap under style drift, that each invert a finding
when uncontrolled. This is the study closest to ours; it is why we separate
finished accuracy from termination (Section~\ref{sec:anatomy}). We differ in
scale (a 1.5B model), in the verifier (a rule-based one whose rejection defect we
measured) and in auditing what that verifier's filtering does to the training
pool, not only whether accuracy moved. Kim et
al.~\cite{ref_selfdistill} show that self-distillation with a teacher that sees the
solution can cut accuracy; our train-on-gold drop resembles this, but we did not test its cause (Section~\ref{sec:final}). Our per-round false-reject measurement
(Section~\ref{sec:audit}) is not new (it is the per-step tracking of
Egashira et al.\ and TinyV, applied to offline SFT and to a defect of our own
pipeline), and we claim neither it nor the patch control as a contribution. Supplementary Appendix~\ref{app:related} gives the full comparison.

\section{Setup}
\label{sec:setup}

\subsection{The GVT Loop}
\label{sec:loop}

Each round $t$ has four steps. \emph{Generate:} the current model samples $k$
solutions per problem. \emph{Verify:} \texttt{math-verify} extracts the final
answer and compares it with the gold. \emph{Select:} accepted solutions are kept.
\emph{Train:} a \emph{fresh} LoRA adapter is attached to the \emph{base} model and
trained on the kept solutions for one epoch. Round~0 is the untrained base model;
the round-$t$ model ($t\ge1$) is the base plus an adapter trained on the pool
accepted at round $t-1$. Adapters are not accumulated, so earlier rounds affect
later ones only through the data they generate; conclusions about
(non-)compounding below apply to this restart-from-base design, which gives
compounding less room than continued training would.

\subsection{The Verifier and Its Directional Bias}
\label{sec:verifier}

A format probe that re-renders gold answers in individually valid forms shows
that \verb|\dfrac|/\verb|\tfrac| are rejected $100\%$ of the time on both
datasets (the command is not parsed). On real solutions (one public
Qwen2.5-7B-Instruct solution per problem for $200$ MATH-500 problems, judged
correct by a scripted normalised comparison) the false-negative rate is
$7.8\%$ ($10$ of $128$; Wilson interval $4.3$--$13.8\%$) on MATH-500 and $0\%$
on GSM8K. Eight of the ten fail when the whole solution is verified but pass
when only the boxed answer is; most are compound answers, so they are probably
largely the gold-storage defect described below rather than extraction of a
wrong token. The audit inspected only rejections, so it gives no
false-positive rate. The verifier also cannot parse $51$ of the $500$ stored
MATH-500 golds ($10.2\%$); under $V$ these problems always count as unsolved and can never enter the
training pool ($V'$ makes $46$ of them parseable for boxed predictions). The offset cancels in within-run differences only if
those problems would not otherwise change, which we did not check, and not in
ratios or levels across partitions ($11$--$18$ such problems per exam partition).

\emph{The patch.} For the control we rewrote \verb|\dfrac| and \verb|\tfrac| to
\verb|\frac| before parsing, removing the $100\%$-rejection class and nothing
else.

\emph{The gold-storage defect (post hoc).} Our pipeline stores the gold without
a delimiter, and \texttt{math-verify} applies its full \LaTeX{} extraction only
inside \verb|\boxed{}| or \verb|$...$|, so a correct boxed prediction can fail
against a bare gold (e.g.\ the prediction \verb|\boxed{\frac{4}{3}}| against
the stored gold \verb|\frac43|). Re-scoring the saved generations with a
verifier $V'$ that wraps the gold in \verb|\boxed{}| when the prediction
contains one shows that this defect rejects $15.5$--$18.3\%$ of the
generations $V'$ accepts. It would enlarge the round-0 kept pool by about a
fifth to a quarter ($440\to542$, $+23\%$; this pool covers only the $350$ selection problems, so its $18.8\%$ rejected share differs from the $15.5$--$18.3\%$ over all $500$), against $0$--$4$ examples for the \verb|\dfrac| patch, and
it lowers absolute pass@1 by $4.7$--$8.0$ points in self-training. It is bound
to the problem, not to the answer: on the $54$--$62$ problems where it ever
fires, it rejects $94$--$98\%$ of correct answers in every round ($87$--$98\%$ where $V$ parses the gold), and $29$--$34$
of these problems have a gold the original verifier cannot parse at all, so no
answer form avoids it there. Re-scored, sample accuracy still falls in all five
runs we re-scored (three self-training, two train-on-gold); the original-verifier changes are $0.73$--$0.85$ of the
re-scored ones, an observed ratio that we do not use to rescale other numbers
or intervals. All tables keep the original verifier, the one the loop trained
with. Section~\ref{sec:audit} uses this defect for the adaptation audit;
Supplementary Appendix~\ref{app:verifier} gives the full analysis.

\subsection{Two Experimental Designs}
\label{sec:designs}

\emph{Design~P (pooled)} loads both datasets ($500$ problems each) together and
trains \emph{one} adapter per round on the union of accepted solutions, with
selection and evaluation on the same $1{,}000$ problems. \emph{Design~I
(isolated)} uses one dataset and one adapter per run, and holds out $30\%$ of the
$500$ problems ($150$) as an \texttt{exam} set; generation runs on all $500$
problems each round, but only the other $350$ can enter the training pool, and
all Design~I pass rates are on the exam set. Table~\ref{tab:designs} summarises
them.

Design~I has three runs, all Qwen2.5-1.5B-Instruct, five rounds, $k=8$, differing
in one respect each: \textbf{Run~A} (MATH-500, original verifier), \textbf{Run~B}
(MATH-500, patched verifier) and \textbf{Run~C} (GSM8K, original verifier). Runs
A and B were repeated at a second seed; since the seed also sets the train/exam
partition, the repetition draws a different 150-problem exam set. Most of the
later evidence comes from higher-rate runs introduced in
Section~\ref{sec:lrcheck}; Table~\ref{tab:inventory} (Supplementary Appendix~\ref{app:isolated}) lists every run. A
``seed'' sets the exam partition, the selection set and the training randomness
at once; runs with the same seed share the partition. ``Five rounds'' means
round~0 (the untrained model) and four adapters, each trained once on the
previous round's pool; ``three rounds'' means round~0 and two adapters, so the
round-2 model is a single adapter (about $195$ optimizer steps for Run~A at
$5\times10^{-4}$) trained on the round-1 pool. All runs use one model family,
Qwen2.5.

\subsection{Models, Data, Training, Metrics}
\label{sec:config}

We use Qwen2.5-0.5B-Instruct and Qwen2.5-1.5B-Instruct (Design~P uses both,
Design~I the 1.5B model only) and the first 500 problems of the GSM8K test
split~\cite{ref_gsm8k} and MATH-500, the 500-problem subset of
MATH~\cite{ref_math}. The prompt is a single user turn in the model's chat
template: ``Solve the math problem. Put the final answer in
\verb|\boxed{}|.'' followed by the problem. Generation samples at temperature
$0.8$ with a cap of $512$ new tokens; the other decoding settings come from the
checkpoint's generation configuration (top-$p$ $0.8$, top-$k$ $20$, repetition
penalty $1.1$). Training is LoRA SFT (\texttt{transformers} 5.0.0, \texttt{peft}
0.19.1, TRL 1.10.0 \texttt{SFTTrainer}; rank $16$, $\alpha=32$, on
\texttt{q\_proj} and \texttt{v\_proj}, LoRA dropout $0$, weight decay $0$, gradient
clipping $1.0$ and trainer seed $42$, all library defaults read from the pinned sources; batch size $2$, no accumulation; one epoch; AdamW with linear learning-rate
decay and no warmup; Python and torch seeds set from the run seed; loss over the whole chat sequence, prompt included; fp16 on
one Kaggle T4), with a training sequence limit of $512$ tokens that includes the
prompt. Because this limit equals the generation cap but also holds the prompt,
any kept solution longer than about $512$ tokens minus the prompt length is cut
in training and loses its final answer and end-of-sequence token
($24$--$26\%$ of the round-0 pool; Section~\ref{sec:anatomy}). All kept
solutions are used, so the number of optimizer steps per round
is half the pool size. The learning rate is $10^{-5}$ unless stated otherwise;
this value was the default of our training code and was not tuned
(Section~\ref{sec:lrcheck}). Metrics per round, on the exam set: pass@1 is the
correctness of the \emph{first} of the $k$ samples (one draw per problem, so it
carries sampling noise); pass@$k$; \emph{sample accuracy}, the mean over all $k$
samples; style counts over all generations (boxed, dfrac, dollar, \dots); and
per-problem added/lost counts between consecutive rounds, where a problem is
\emph{solved} in a round if at least one of its $k$ samples is accepted. Added
minus lost therefore equals $150\times\Delta$pass@$k$, so these counts add only
the gross churn behind the net change, not an independent test. pass@1 was
fixed in advance; sample accuracy was not among the pre-set measures, but it
estimates the same quantity from all $k$ samples with less variance, so after
the runs we chose it as the primary accuracy measure. Claims that rest on it,
including the $2.2$-point bound, are therefore post hoc in the choice of measure. ``Correct'' always means accepted by the verifier (preset
\texttt{reward}), the same verifier that selects the training data.


\paragraph{Uncertainty.} For every within-run change (round~0 against the last
round) and every paired comparison of two runs on the same exam partition we
report a 95\% interval from a paired problem-level bootstrap ($10{,}000$
resamples of the $150$ exam problems; $2{,}000$ for the matched and audit columns), which respects the clustering of the $k$
samples within a problem. These intervals do \emph{not} cover training-seed
variance, which we see only across two or three seeds: the replicate
pair we use as a gauge ends $3.2$ points of sample accuracy apart (a seed also changes the exam
partition, so such pairs mix training and partition variance). We make no multiplicity correction and the runs are not independent (they share
exam partitions and round-0 draws), so patterns across runs are descriptive, not tests;
we apply the same reading to declines and to gains. Single proportions use Wilson
intervals. \emph{Detectable effects.} From the interval half-widths $h$ in
Table~\ref{tab:ci}, a change would be detected with $80\%$ power at the 5\%
level only if it exceeded about $2.8\,h/1.96$: about $10$ points of pass@1,
$8$--$10$ of pass@$k$ and $3.5$--$5.5$ of sample accuracy on MATH-500. Because
$h$ omits training-seed variance (the replicate pair differs by $3.2$ points and
B minus A flips from $+5.3$ to $-7.3$ on near-identical pools,
Section~\ref{sec:lrcheck}), these are optimistic thresholds, not bounds: even with $h$ as the whole error, a gain
at that size is missed $20\%$ of the time, and probably more often (Supplementary Appendix~\ref{app:uncertainty}).

\section{Results}
\label{sec:results}

\subsection{Design P: An Apparent Collapse}
\label{sec:pooled}

In the pooled design the 1.5B model declines severely and with acceleration
(pass@1 $50.0\%\to18.6\%$, pass@8 $72.6\%\to48.6\%$; Table~\ref{tab:passk}, Supplementary Appendix~\ref{app:pooled}),
on both datasets (Table~\ref{tab:bydataset}); the 0.5B decline is small. The
accepted pool contracts every round and \verb|\boxed{}| falls from $5{,}608$ to
$2{,}693$ of $8{,}000$ generations. This design cannot be read causally: one
adapter trains on both datasets' pools, it scores on the problems it selects
from, and it trains on about $4.5\times$ as many examples per round as the
isolated MATH-500 run ($3{,}962$ against $872$ at round~0); no run separates these
(Supplementary Appendix~\ref{app:pooled}).

\subsection{Design I at $10^{-5}$: Small Declines, a Larger One on GSM8K}
\label{sec:isolated}

On the held-out exam (Fig.~\ref{fig:isolated} and Table~\ref{tab:isolated} in Supplementary Appendix~\ref{app:isolated}) pass@8
ends within $4$ points of round~0 in every run. On MATH-500, pass@1 is
$2.7$--$4.0$ points lower at round~4, with every interval including zero
(Table~\ref{tab:ci}); sample accuracy falls $1.2$--$3.2$ points, with an
interval excluding zero only for B at seed~0 (A at seed~0 reaches $0.0$), and
matched finished accuracy $1.1$--$4.2$, with intervals that only just exclude
zero for A and B at seed~0, which are close to one replicate. We read this as a
possible small decline, not an established one; what we do not see is a gain
of the size the design could detect. Run~C
(GSM8K) declines more: pass@1 falls $12.7$ points ($[-21.3,-4.0]$), almost all in
finished accuracy, but it took about twice Run~A's optimizer steps, so we do not
attribute this to the dataset alone. Length and cap-hit share barely change on MATH-500, weak evidence that the adapter
did little (Supplementary Appendix~\ref{app:isolated}).

\subsection{The Patch Control Could Not Detect an Effect}
\label{sec:control}
\label{sec:prevalence}
\label{sec:answer-diversity}

Runs A and B differ only in the \verb|\dfrac|/\verb|\tfrac| patch. At round~4,
B minus A is $-0.7$ [$-3.3$, $+2.0$] points of pass@1 at seed~0 and $+0.0$
[$-2.7$, $+2.7$] at seed~1 (recomputed from the saved generations; Table~\ref{tab:ci} lists run-log values, $-0.7$ at seed~1). This null is uninformative: \verb|\dfrac| appears
in $0.50$--$0.73\%$ of generations, so the kept pools differ by $3$--$4$
examples (under $0.5\%$ of the $872$-example pool), and at $10^{-5}$ the adapter changes the outputs
little anyway, so a null is what prevalence alone predicts; the frequent gold-storage defect is in both arms.

\subsection{Higher Learning Rates}
\label{sec:lrcheck}

Because $10^{-5}$ changed the MATH-500 outputs little, we reran Run~A at
$10^{-4}$ (three seeds), $2\times10^{-4}$ (two) and $5\times10^{-4}$ (two,
with Run~B), with three rounds and $k=4$ (Table~\ref{tab:lrseries}). These
runs differ from the $10^{-5}$ runs in $k$, rounds, pool size and the pass@$k$
metric as well as rate, so only $10^{-4}$ to $5\times10^{-4}$ isolates the rate.
\emph{At $5\times10^{-4}$ every A/B run declines:} pass@1 falls $4.0$--$11.3$
points and the sample-accuracy intervals exclude zero in all four A/B runs,
though same-seed A and B are closer to two replicates than four.
\emph{The decline may grow with the rate on one seed} (sample-accuracy changes at $10^{-5}$ and $10^{-4}$ are similar, $-1.2$ to $-3.2$ and $-1.0$ to $-3.7$): at seed~1 sample accuracy
falls $1.5$, $3.5$ and $8.3$ points across the three rates, at seed~2 an almost
flat $3.7$, $3.8$ and $4.7$. The cap-hit share rises only at $5\times10^{-4}$.
\emph{The patch has no consistent effect:} B minus A in round-2 pass@1 is
$+5.3$ at seed~1 and $-7.3$ [$-14.0$, $-0.7$] at seed~2, a sign flip from pools
differing by $0$--$3$ examples: trajectory noise.
\emph{Style does not separate imitation from drift:} boxed share falls with the
pool, which both would produce, and the no-filter and gold arms decline with no verifier in the selection step (they are still scored by it), so generic SFT damage fits as well.

\begin{table}[!htb]
\caption{Run~A (MATH-500, original verifier, $k=4$, three rounds) across learning
rates and seeds: pass@1, pass@4 and sample accuracy (600 exam generations) at
round~0 $\to$ round~2, the share of round-2 generations that hit the
$512$-token cap (decoded text re-tokenised to $\ge508$ tokens; round~0 is
$43.0$--$43.2\%$, $48.2\%$ at seed~3), and problems added/lost summed over
rounds 1--2. The two $5\times10^{-4}$ rows are the A rows
of Table~\ref{tab:lrcheck}.}
\label{tab:lrseries}
\centering
\setlength{\tabcolsep}{3pt}
\begin{tabular}{llccccc}
\toprule
Rate & Seed & pass@1 & pass@4 & Sample acc. & Cap-hit (R2) & Added/lost \\
\midrule
$10^{-4}$ & 1 & $33.3\to34.0$ & $50.0\to49.3$ & $33.7\to32.2$ & $42.7\%$ & $20/21$ \\
$10^{-4}$ & 2 & $31.3\to28.0$ & $51.3\to45.3$ & $30.8\to27.2$ & $42.8\%$ & $25/34$ \\
$10^{-4}$ & 3 & $29.3\to26.7$ & $40.0\to41.3$ & $28.5\to27.5$ & $45.7\%$ & $16/14$ \\
$2\times10^{-4}$ & 1 & $33.3\to31.3$ & $50.0\to46.7$ & $33.7\to30.2$ & $40.8\%$ & $23/28$ \\
$2\times10^{-4}$ & 2 & $31.3\to25.3$ & $51.3\to45.3$ & $30.8\to27.0$ & $39.7\%$ & $21/30$ \\
$5\times10^{-4}$ & 1 & $33.3\to23.3$ & $50.0\to42.7$ & $33.7\to25.3$ & $46.2\%$ & $22/33$ \\
$5\times10^{-4}$ & 2 & $31.3\to27.3$ & $51.3\to44.7$ & $30.8\to26.2$ & $51.2\%$ & $26/36$ \\
\bottomrule
\end{tabular}
\end{table}

\begin{table}[!htb]
\caption{Change from round~0 to the last round (percentage points) with 95\%
paired problem-level bootstrap intervals over the 150 exam problems
($10{,}000$ resamples; \texttt{bootstrap\_ci.py} in the code repository). Last round
is R4 at $10^{-5}$ ($k=8$) and R2 elsewhere ($k=4$). ``Finished, matched'' is
the last-round accuracy of generations below the $512$-token cap minus the
untrained model's accuracy on the same problem mix (Section~\ref{sec:anatomy};
$2{,}000$ resamples, \texttt{capanat.py} in the repository). The
intervals do not include training-seed variance. Recomputed values match the
reported ones except six round values that differ by one problem
(Supplementary Appendix~\ref{app:uncertainty}), e.g.\ B, seed~1 at $10^{-5}$, whose round-4 pass@8
is $58.0\%$ against $58.7\%$; the matched column uses the separate re-scoring
described there.}
\label{tab:ci}
\centering
\scriptsize
\setlength{\tabcolsep}{1.5pt}
\begin{tabular}{lcccc}
\toprule
Run, seed & $\Delta$pass@1 & $\Delta$pass@$k$ & $\Delta$sample acc. & Finished, matched \\
\midrule
\multicolumn{5}{l}{\emph{Learning rate $10^{-5}$}} \\
A, 0 & $-2.7$ [$-9.3,4.0$] & $-1.3$ [$-8.0,5.3$] & $-2.4$ [$-4.8,0.0$] & $-3.8$ [$-7.2,-0.5$] \\
A, 1 & $-3.3$ [$-10.0,3.3$] & $+0.0$ [$-6.0,6.0$] & $-1.5$ [$-4.1,1.0$] & $-2.3$ [$-5.9,1.2$] \\
B, 0 & $-3.3$ [$-10.0,3.3$] & $-4.0$ [$-10.7,2.7$] & $-3.2$ [$-5.9,-0.6$] & $-4.2$ [$-7.8,-0.9$] \\
B, 1 & $-4.0$ [$-10.7,2.7$] & $-1.3$ [$-7.3,4.7$] & $-1.2$ [$-3.7,1.3$] & $-1.1$ [$-4.6,2.6$] \\
C, 0 & $-12.7$ [$-21.3,-4.0$] & $-2.7$ [$-7.3,2.0$] & $-13.1$ [$-16.7,-9.8$] & $-12.8$ [$-16.5,-9.4$] \\
\midrule
\multicolumn{5}{l}{\emph{Learning rate $10^{-4}$}} \\
A, 1 & $+0.7$ [$-6.7,8.0$] & $-0.7$ [$-6.0,4.7$] & $-1.5$ [$-5.3,2.2$] & $-4.3$ [$-9.6,1.1$] \\
A, 2 & $-3.3$ [$-9.3,2.7$] & $-6.0$ [$-12.7,0.7$] & $-3.7$ [$-7.3,-0.2$] & $-3.5$ [$-9.2,2.3$] \\
A, 3 & $-2.7$ [$-10.0,4.0$] & $+1.3$ [$-4.7,7.3$] & $-1.0$ [$-4.3,2.2$] & $-0.6$ [$-6.0,4.7$] \\
\multicolumn{5}{l}{\emph{Learning rate $2\!\times\!10^{-4}$}} \\
A, 1 & $-2.0$ [$-8.0,4.0$] & $-3.3$ [$-10.0,2.7$] & $-3.5$ [$-7.2,0.0$] & $-4.0$ [$-8.9,0.5$] \\
A, 2 & $-6.0$ [$-12.7,0.7$] & $-6.0$ [$-12.7,0.7$] & $-3.8$ [$-7.5,-0.2$] & $-4.7$ [$-10.1,1.0$] \\
\midrule
\multicolumn{5}{l}{\emph{Learning rate $5\!\times\!10^{-4}$}} \\
A, 1 & $-10.0$ [$-16.7,-3.3$] & $-7.3$ [$-14.0,0.0$] & $-8.3$ [$-12.3,-4.5$] & $-9.1$ [$-14.2,-4.2$] \\
A, 2 & $-4.0$ [$-11.3,3.3$] & $-6.7$ [$-13.3,0.0$] & $-4.7$ [$-8.2,-1.2$] & $-5.5$ [$-11.7,0.1$] \\
B, 1 & $-4.7$ [$-12.0,2.7$] & $-2.0$ [$-8.7,4.7$] & $-6.7$ [$-10.5,-2.8$] & $-6.8$ [$-12.4,-1.0$] \\
B, 2 & $-11.3$ [$-17.3,-5.3$] & $-11.3$ [$-18.0,-4.7$] & $-9.2$ [$-12.8,-5.5$] & $-5.9$ [$-10.9,-1.0$] \\
No filter, 1 & $-6.0$ [$-12.7,1.3$] & $-9.3$ [$-16.7,-2.7$] & $-8.2$ [$-12.2,-4.3$] & $-10.6$ [$-16.4,-5.5$] \\
No filter, 2 & $-8.7$ [$-16.0,-1.3$] & $-12.0$ [$-19.3,-4.7$] & $-10.2$ [$-14.2,-6.3$] & $-6.6$ [$-11.4,-2.0$] \\
Gold, 1 & $-18.0$ [$-26.7,-10.0$] & $-20.7$ [$-28.7,-12.7$] & $-19.7$ [$-25.3,-14.2$] & $-25.6$ [$-32.1,-18.9$] \\
Gold, 2 & $-21.3$ [$-28.7,-14.0$] & $-20.7$ [$-29.3,-12.7$] & $-18.7$ [$-23.7,-13.7$] & $-25.9$ [$-31.7,-20.0$] \\
C, 1 & $-22.0$ [$-29.3,-15.3$] & $-9.3$ [$-15.3,-3.3$] & $-15.3$ [$-19.8,-11.2$] & $-12.1$ [$-16.6,-7.9$] \\
\bottomrule
\end{tabular}
\end{table}

\paragraph{Two controls at $5\times10^{-4}$.}
The \emph{no-filter} run trains on all $k$ samples of every selection problem
($1{,}400$ examples per round), removing the verifier; the
\emph{train-on-gold} run trains on the $350$ reference solutions every round
(Table~\ref{tab:controls}). Neither is a single-variable control: no-filter
takes $3$--$4\times$ Run~A's optimizer steps, train-on-gold about the same.
Removing the filter is not better than keeping it beyond noise and is lower on
pass@4 at both seeds (round-2 differences $-2.0$ and $-5.3$ points, intervals
including zero), but the dose difference means this cannot be attributed to the
filter; the no-filter pool is also truncated far more often in training
($62.4\%$ against $24.5\%$ at round~0). Training on reference solutions is
followed by a much larger decline (pass@4 down about $21$ points at both
seeds) while the model adopts the reference style (\texttt{\$}-delimited
spans $5\%\to98\%$). Post hoc checks bound verifier error from below only
(about $2$ points of sample accuracy among boxed answers), and the $14\%$ of
round-2 generations without \verb|\boxed| are not covered, so we think this
fall is unlikely to be mainly verifier error but cannot bound it from above.
That no arm, this one included, improves on the untrained model means we have
no positive reference: the nulls above may reflect the recipe (LoRA on
$q$/$v$ projections, loss over the prompt, a $512$-token limit) rather than
self-training. Run~C at
$5\times10^{-4}$ (GSM8K, one run) falls $22.0$ points of pass@1, from both less
accurate and more unfinished solutions (Supplementary Appendix~\ref{app:lrcheck}).

\begin{table}[!htb]
\caption{Controls at $5\times10^{-4}$ (MATH-500, $k=4$, three rounds), on the
exam partition of Run~A at the same seed. Pool is the kept pool at round~0 $\to$ round~2; the round-2 model trains on the
round-1 pool ($392$ and $386$ examples, about $195$ steps) and optimizer steps are
half the pool size.}
\label{tab:controls}
\centering
\setlength{\tabcolsep}{5pt}
\begin{tabular}{lccccccc}
\toprule
Run & Pool & \multicolumn{3}{c}{pass@1} & \multicolumn{3}{c}{pass@4} \\
 &  & R0 & R1 & R2 & R0 & R1 & R2 \\
\midrule
\multicolumn{8}{l}{\emph{Seed 1}} \\
A (verifier-filtered) & 440$\to$326 & 33.3 & 32.7 & 23.3 & 50.0 & 49.3 & 42.7 \\
No filter             & 1400 & 33.3 & 26.7 & 27.3 & 50.0 & 49.3 & 40.7 \\
Train on gold         & 350  & 33.3 & 18.0 & 15.3 & 50.0 & 33.3 & 29.3 \\
\midrule
\multicolumn{8}{l}{\emph{Seed 2}} \\
A (verifier-filtered) & 424$\to$350 & 31.3 & 31.3 & 27.3 & 51.3 & 46.0 & 44.7 \\
No filter             & 1400 & 31.3 & 34.7 & 22.7 & 51.3 & 48.7 & 39.3 \\
Train on gold         & 350  & 31.3 & 8.7 & 10.0 & 51.3 & 28.7 & 30.7 \\
\bottomrule
\end{tabular}
\end{table}

\subsection{Anatomy of the Decline}
\label{sec:anatomy}

This subsection is post hoc and descriptive. Sample accuracy mixes whether a
finished solution is right with whether the model finishes within the
$512$-token cap at all. Re-tokenising the saved text, we call a generation
\emph{cap-hit} at $\ge508$ tokens. At the rates where the share does not rise, $40$--$46\%$ of round-2 MATH-500 generations
already hit the cap (Table~\ref{tab:lrseries}), so the cap binds from the start. At $5\times10^{-4}$, compared with the untrained model on the same
problem mix, finished generations lose $5.5$--$10.6$ points, with intervals
excluding zero in five of six runs (the six include two same-seed A/B pairs,
which are close to replicates), while the cap-hit share rises $2$--$14$ points.
The matched measure conditions on which problems finish after training, so it is
descriptive. A three-part decomposition attributes $12$--$30\%$ of the decline to the cap share in the
three original seed-1 runs (and $-2$ and $5\%$ in two further runs) and $52$--$72\%$ at seed~2 (Table~\ref{tab:decomp}, Supplementary Appendix~\ref{app:anatomy}). At
$10^{-4}$ and $2\times10^{-4}$ the cap-hit share does not rise and finished
accuracy carries the whole change. Train-on-gold is different: responses
shorten by more than half and finished accuracy collapses. Retraining without truncated examples, or re-scoring at a $1{,}536$-token
cap, leaves the deficit in place (Supplementary Appendix~\ref{app:anatomy}).

\subsection{A Final Round with a Revised Recipe}
\label{sec:final}
Because no earlier arm improved, we ran one last round with several recipe changes made at once,
aimed at the suspected causes (no ablation of any single one): loss on the completion only, LoRA on all linear layers,
a 1,024-token cap and training limit, normalised gold answers, and a pool of
300 problems from the MATH training split, so that the exam is all 500
MATH-500 problems. It is $10^{-4}$ with $k=2$ and a single adapter (round~0 to round~1); the loop has three seeds, the gate one. We report no overlap check of the $300$ pool problems against MATH-500. The
verifier-filtered loop (298 examples kept) is compared with training on the 300
reference solutions, a positive gate that must raise accuracy for the recipe to
count as able to help. 

\begin{table}[!htb]
\caption{Final round, MATH-500 (500 exam problems, $k=2$, $10^{-4}$).
Change from each run's own untrained round~0 to round 1 in points (round-0 pass@1 is
$49.2$, $46.8$, $48.2$ at seeds 1--3, from sampling), with 95\% paired
problem-level bootstrap intervals that omit training-seed variance. The gold arm has one seed.}
\label{tab:final}
\centering
\resizebox{\linewidth}{!}{%
\begin{tabular}{lccc}
\hline
Arm & pass@1 & pass@2 & Sample acc.\\
\hline
Untrained (seed 1) & 49.2 & 59.4 & --\\
Loop, seed 1 & $-4.8$ $[-9.0,-0.6]$ & $-3.0$ $[-6.8,+0.6]$ & $-2.9$ $[-5.8,0.0]$\\
Loop, seed 2 & $-3.2$ $[-6.8,+0.6]$ & $-6.2$ $[-9.8,-2.4]$ & $-5.3$ $[-8.1,-2.4]$\\
Loop, seed 3 & $-4.2$ $[-8.4,-0.2]$ & $-5.6$ $[-9.4,-1.6]$ & $-5.4$ $[-8.3,-2.4]$\\
Loop, mean of seeds & $-4.1$ & $-4.9$ & $-4.5$\\
Train on gold & $-32.2$ $[-36.8,-27.6]$ & $-32.0$ $[-36.6,-27.2]$ & $-30.2$ $[-33.8,-26.6]$\\
\hline
\end{tabular}}
\end{table}

The gate fails. Training on reference solutions lowers pass@1 from $49.2\%$ to
$17.0\%$, and the model adopts the reference style ($\$$-delimited spans in
$60$ of $1{,}600$ generations at round 0, $1{,}577$ at round 1, counted over the
exam and selection problems), as in the earlier train-on-gold arms. A
rerun at $2\times10^{-5}$ fails the same way ($49.2\%\to20.6\%$, $\$$-spans in $1{,}558$ of
$1{,}600$ generations), so the failure is not specific to the rate. We noticed the shift to the reference
style but did not test it as the cause. That style carries no \verb|\boxed| answer, which both graders read, and a collapse of this size from $300$ reference
solutions is also what a pipeline defect would produce. The gate therefore fails to validate the
revised recipe, and the recipe cannot be called corrected. We have no recipe known to
raise accuracy, so the loop's decline cannot be attributed to self-training or to the verifier. The verifier-filtered
loop loses pass@1 in all three seeds ($4.8$, $3.2$, $4.2$ points; two intervals exclude zero) and
sample accuracy in all three (intervals exclude zero for seeds 2 and 3). The mean pass@1 change over the three seeds is $-4.1$ ($t$ interval
$[-6.1,-2.1]$), and $-4.5$ for sample accuracy ($[-8.1,-1.0]$). The loop does not
help; the point estimates are negative in all three seeds, but the sample is three seeds with one
exam draw and one rate, so we do not claim that it costs accuracy, and we did not isolate the cause. The detectable effect
is about $6$ points of pass@1 ($2.8h/1.96$), against about $10$ before.

\emph{Injected bias.} As intended positive controls for the audit (neither turns out to be one) we reran the loop once
each (seed~1, $10^{-4}$) with a selection verifier that rejects \verb|\boxed|
answers (exam verifier clean). In the first, a hash-chosen $30\%$ of problems are
affected: the \verb|\boxed| share falls $4.9$ points on them and $3.1$ on the others
(difference $-1.8$, $[-7.4,3.8]$); a hash cannot be learned from the problem, so this
control could not elicit selective avoidance. In the second the rule is visible: every
\verb|\boxed| answer is rejected when the gold contains \verb|\frac| ($49$ of the $300$ selection
problems, $263$ examples kept instead of $298$; $70$ of the $500$ exam problems, where the shares below are measured). The share on those problems stays at
$95.0\%$ and on the others goes from $90.2\%$ to $87.4\%$ (difference $+2.8$,
$[-2.8,+8.3]$), against $-3.7$, $-2.7$ and $-5.5$ for the same split in the three unbiased
seeds, a shift in the direction opposite to avoidance; pass@1 falls $7.0$ points ($[-11.2,-3.0]$),
more than in the unbiased seeds, which we do not interpret (one run). In neither run does the
\verb|\boxed| share fall selectively where boxing is penalised, so neither
validates the audit.


\section{Discussion}
\label{sec:discussion}

\subsection{Answer to the Research Question}
\label{sec:answer}

\emph{Does the model get better at mathematics?} Not detectably. No run raises
pass@1 or pass@$k$ with an interval excluding zero, at any rate; in the
500-problem final round pass@1 falls $3.2$--$4.8$ points in each of three seeds. The sample-accuracy figure of
about $2.2$ points is the largest upper limit among observed intervals, a post hoc
figure that omits seed variance, not a power statement; on the pre-set measures the design could detect only about $10$ points of pass@1 and
$8$--$10$ of pass@$k$ ($6$ and $5$ in the final round). Because no arm
improves and the positive gate failed, the null belongs to this recipe, and the
study cannot say whether self-training in general helps. Every score except the final-round check below is graded by the verifier that selects
the data (re-scoring five runs, three self-training and two train-on-gold, with $V'$ leaves the decline in place, Section~\ref{sec:setup}). On the final-round generations a
second grader (normalised string match of the last \verb|\boxed| against the gold, no \texttt{math-verify};
\texttt{indep\_grade.py}) gives pass@1 changes of $-3.6$, $-2.2$, $-4.6$ for the three seeds ($-30.0$ for the gold arm) and
sample-accuracy changes of $-2.5$, $-5.5$, $-5.3$ (same sign and size). It reads only the last \verb|\boxed|, so a falling boxed share lowers both graders, and it does not rule out a format artefact in the gate. Earlier runs are not independently graded.

\emph{Does it learn to write what the verifier accepts?} The data cannot settle
this. The pre-set pattern requires pass@1 to rise, which never happens
detectably (largest rise $+0.7$ [$-6.7$, $+8.0$]), so its silence is no evidence either way;
no run varies the frequent gold-storage defect, and the audit of Section~\ref{sec:audit} is too weak.

\subsection{Interpretation}
\label{sec:threevars}

One explanation of the Design~P curves is a self-reinforcing
shrinking pool shaped by the biased verifier. The data neither support nor rule
it out: the patch tests a rare bias, and the no-filter
control declines without any selection verifier but with more dose (and, with a constant pool, cannot test the shrinking pool). At this resolution, \emph{removing a rare format defect, or the
filter entirely, made no consistent difference} (no filter is lower on pass@$k$ at $3$--$4\times$ the dose); neither tests the systematic
false positives Egashira et al.\ find harmful in RLVR.

\subsection{A Proposed Audit for GVT with a Biased Verifier}
\label{sec:audit}

Our pre-set rule tied ``learning the verifier's style'' to rising accuracy, so
it could not speak once accuracy fell. From post hoc analyses we propose four
checks for self-training with a rule-based verifier, as a protocol for future
studies, not as tested methods.

\emph{(i) A per-round false-reject audit.} Fix a repaired verifier $V'$ that
corrects a known defect class of the training verifier $V$. At each round $t$,
on the loop's own generations, compute
\[
r_t = \frac{\#\{\text{accepted by } V' \text{ but rejected by } V\}}{\#\{\text{accepted by } V'\}} .
\]
If the model learns to write what $V$ accepts, $r_t$ may fall over rounds, but
it also falls when fewer answers are boxed: since $V'$ differs from $V$ only on boxed
predictions, $r_t=r_t^{\text{boxed}}\times(\text{boxed share among }V'\text{-accepted})$ holds as an identity. This is the per-step tracking of
RLVR studies~\cite{ref_delayplateau,ref_tinyv} applied to offline SFT (Table~\ref{tab:rt}; \texttt{rt\_audit.py}).
For the gold-storage defect it falls in all three self-training runs we re-scored,
by $2.6$, $2.6$ and $1.5$ points, with intervals reaching or just touching
zero. The fall is accounted for by the falling
\verb|\boxed| share among correct answers ($91\to76\%$ and $93\to78\%$ at
$5\times10^{-4}$): an identity, not a test, which does not exclude
adaptation. On boxed answers alone the changes ($-0.4$ to $+0.7$ points) have intervals
($\approx\pm3$) that contain the overall fall of the same run ($-2.6$, $-2.6$, $-1.5$), so ``flat'' means not detected. The defect is bound to the problem: it fires on $54$--$62$
problems, rejecting $94$--$98\%$ of correct answers, and $29$--$34$ of them have an unparseable gold, so no answer
form escapes it. On the $24$--$28$ affected
problems where some form could, $r_t$ moves by $-5.5$, $-1.6$ and $-0.2$
points (intervals $10$--$26$ wide), and the pool
holds only $2$--$12$ such unboxed examples per round, so we expect weak pressure to adapt (untested; no dose--response). The two injected-bias runs
(Section~\ref{sec:final}) show no selective avoidance, but neither can tell us whether the audit would detect it.
Train-on-gold, with no verifier in the loop, gives
$r_t$ rises of $2.1$--$3.3$ points ($2.8$ and $1.7$ on boxed answers alone;
Table~\ref{tab:rt}), so $r_t$ drifts as much as the falls we observe, unrelated to
the verifier. It should be reported on
the answers and problems where $V$ and $V'$ can differ and an accepted form
exists, with the count of training examples that could teach avoidance.

\begin{table}[!htb]
\caption{The false-reject audit for the gold-storage defect (post hoc), in
percent (all $500$ problems, selection problems included), round~0 $\to$ last round, with the change and a $95\%$ problem-level bootstrap interval ($2{,}000$
resamples). ``Boxed only'' restricts to predictions containing
\texttt{\textbackslash boxed}; ``affected, parseable'' to problems where the
defect ever fires and $V$ can parse the gold. Pool: examples per round from
affected problems that $V$ accepts (and so avoid the defect).}
\label{tab:rt}
\centering
\footnotesize
\setlength{\tabcolsep}{4pt}
\begin{tabular}{llrrrrc}
\toprule
Run & $r_t$ subset & R0 & Last & $\Delta$ & 95\% interval & Pool \\
\midrule
A, $5\times10^{-4}$, seed 1 & all & $18.3$ & $15.7$ & $-2.6$ & $[-5.3, 0.3]$ & $3$--$5$ \\
 & boxed only & $20.0$ & $20.7$ & $+0.7$ & $[-2.9, 4.3]$ & \\
 & affected, parseable & $92.8$ & $87.2$ & $-5.5$ & $[-20.1, 6.3]$ & \\
\addlinespace
A, $5\times10^{-4}$, seed 2 & all & $18.1$ & $15.5$ & $-2.6$ & $[-5.2, 0.1]$ & $2$--$3$ \\
 & boxed only & $19.4$ & $19.8$ & $+0.4$ & $[-2.8, 3.4]$ & \\
 & affected, parseable & $95.1$ & $93.5$ & $-1.6$ & $[-14.7, 9.5]$ & \\
\addlinespace
A, $10^{-5}$, seed 0 & all & $18.3$ & $16.8$ & $-1.5$ & $[-2.9, 0.0]$ & $9$--$12$ \\
 & boxed only & $19.9$ & $19.5$ & $-0.4$ & $[-2.0, 1.2]$ & \\
 & affected, parseable & $92.3$ & $92.0$ & $-0.2$ & $[-4.9, 4.6]$ & \\
\addlinespace
Gold, seed 1 & all & $18.2$ & $21.5$ & $+3.3$ & $[-1.2, 8.1]$ & --- \\
 & boxed only & $20.0$ & $22.8$ & $+2.8$ & $[-2.0, 7.7]$ & \\
 & affected, parseable & $95.5$ & $94.7$ & $-0.7$ & $[-9.1, 6.2]$ & \\
\addlinespace
Gold, seed 2 & all & $18.0$ & $20.2$ & $+2.1$ & $[-2.2, 6.7]$ & --- \\
 & boxed only & $19.4$ & $21.1$ & $+1.7$ & $[-3.3, 6.5]$ & \\
 & affected, parseable & $95.1$ & $97.6$ & $+2.5$ & $[-4.8, 11.9]$ & \\
\bottomrule
\end{tabular}
\end{table}

\emph{(ii) Decompose declines under a fixed cap.} Report the change in sample
accuracy split into a finished-accuracy part and a cap-share part
(Table~\ref{tab:decomp}, Supplementary Appendix~\ref{app:anatomy}). Here the split is descriptive and its balance
differs by partition: at $5\times10^{-4}$ the cap share accounts for
$12$--$30\%$ of the decline in the three original seed-1 runs and $52$--$72\%$
at seed~2.

\emph{(iii) Check the dose of a control first} (our \verb|\dfrac| patch changed $0$--$4$ examples; the unvaried defect a fifth to a quarter of the pool) and \emph{(iv) match dose across designs}, which the pooled design did not.

\section{Limitations}
\label{sec:limitations}

\emph{Unequal training doses} confound every cross-arm comparison, and the
controls are underpowered: the patch changes $0$--$4$ training examples, and
the no-filter and train-on-gold controls have two seeds each. \emph{Intervals
cover only problem draw and sampling,} not training-seed variance; the final-round loop has three seeds, the
gate and the bias controls one, and no multiplicity correction is applied. \emph{The
learning-rate grid is coarse} and changes $k$, rounds and dose with the rate. \emph{Power is limited} (about $10$ points of pass@1 with $150$ exam
problems, about $6$ in Section~\ref{sec:final}). \emph{The adaptation audit is weak:} $24$--$28$
problems, $2$--$12$ avoidance examples per round, and no valid positive control.
\emph{The final-round loop has no no-filter arm.} \emph{The verifier audit} used other models' solutions and inspected only
rejections; the gold-storage defect ran unvaried through every earlier run, whose headline numbers are graded
by the biased verifier that also selects the data, and the final-round string-match grader shares the \verb|\boxed| dependence. The final-round decline and the gate failure cannot be attributed to any single recipe change. \emph{Truncation:} the $512$-token cap binds from
the start (a quarter of the pool is cut). \emph{Scope:}
one model family and size, all instruction-tuned (no base model), one verifier, a
restart-from-base design that leaves compounding little room (no continued
training), possible Qwen2.5 contamination on MATH/GSM8K (unquantified), and a default repetition penalty ($1.1$) that may depress the baseline.

\section{Conclusion}
\label{sec:conclusion}

We asked whether generate--verify--select--train iterations make a 1.5B model
better at mathematics or only teach it to write what a rule-based verifier
accepts. The design detected no pass@1 or pass@$k$ gain at any rate (it could see only gains of about
$10$ and $8$--$10$ points; $6$ in a final $500$-problem round, where pass@1 falls $3.2$--$4.8$ points in
all three seeds). Because training on reference solutions lowers pass@1 by about $30$ points even with a
revised recipe, this is a null for the recipe and says little about self-training, and since scores are
graded by the selecting verifier (apart from one string-match check), neither branch of the question is answered. The post hoc false-reject audit is uninformative about
adaptation: its fall is tied to the falling \verb|\boxed| share by an identity, and it has little power (about $25$
affected problems, a handful of avoidance examples per round, no valid positive control). At the highest rate the model is less accurate on the
solutions it finishes. The pooled design's decline is absent in the isolated runs at $10^{-5}$, which use about $4.5\times$ fewer examples, so dose and
pooling are confounded.

A dose-matched control for the gold-storage defect, a recipe that passes the
positive gate, and a second grader come first (Supplementary Appendix~\ref{app:future}).

\begin{credits}
\subsubsection{\ackname}
This work was carried out as an undergraduate research project. No external
funding was received.

\subsubsection{\discintname}
The authors have no competing interests to declare that are relevant to the
content of this article.

\subsubsection{Code availability}
Code, saved generations and analysis scripts are at
\url{https://github.com/tsan2711/DACNTT_2026} (tag \texttt{paper-2026-11b}). Generation-time token counts, Design~P's training logs and the
$1{,}536$-token re-scoring outputs were not saved.

\subsubsection{Use of AI tools}
The authors used Claude (Anthropic), through Claude Code, to help write the
manuscript, the analysis and training scripts, and to run automated reviews. The authors designed the study, ran the experiments
and take responsibility for the content.
\end{credits}

\begin{thebibliography}{99}

\bibitem{ref_imperfect}
Cai, X.-Q., Wang, W., Liu, F., Liu, T., Niu, G., Sugiyama, M.:
Reinforcement Learning with Verifiable yet Noisy Rewards under Imperfect
Verifiers. arXiv:2510.00915 (2025)

\bibitem{ref_gsm8k}
Cobbe, K., Kosaraju, V., Bavarian, M., Chen, M., Jun, H., Kaiser, L.,
Plappert, M., Tworek, J., Hilton, J., Nakano, R., Hesse, C., Schulman, J.:
Training Verifiers to Solve Math Word Problems. arXiv:2110.14168 (2021)

\bibitem{ref_openrs}
Dang, Q.-A., Ngo, C.: Reinforcement Learning for Reasoning in Small LLMs: What
Works and What Doesn't. arXiv:2503.16219 (2025)

\bibitem{ref_delayplateau}
Egashira, K., Vero, M., Dekoninck, J., D\"orner, F.E., Staab, R., Vechev,
M.: Delay, Plateau, or Collapse: Evaluating the Impact of Systematic
Verification Error on RLVR. arXiv:2605.02909 (2026)

\bibitem{ref_math}
Hendrycks, D., Burns, C., Kadavath, S., Arora, A., Basart, S., Tang, E.,
Song, D., Steinhardt, J.: Measuring Mathematical Problem Solving With the
MATH Dataset. arXiv:2103.03874 (2021)

\bibitem{ref_robustness}
Huang, Y., Zeng, W., Zeng, X., Zhu, Q., He, J.: From Accuracy to
Robustness: A Study of Rule- and Model-based Verifiers in Mathematical
Reasoning. arXiv:2505.22203 (2025)

\bibitem{ref_selfdistill}
Kim, J., Luo, X., Kim, M., Lee, S., Kim, D., Jeon, J., Li, D., Yang, Y.:
Why Does Self-Distillation (Sometimes) Degrade the Reasoning Capability of
LLMs? arXiv:2603.24472 (2026)

\bibitem{ref_selectionbias}
Qiao, X., Du, X., Liu, W., Zhang, J., Mai, P., Zhang, M., Pang, Y.:
When Sample Selection Bias Precipitates Model Collapse. arXiv:2606.13732
(2026)

\bibitem{ref_spurious}
Shao, R., Li, S.S., Xin, R., Geng, S., Wang, Y., Oh, S., Du, S.S., Lambert,
N., Min, S., Krishna, R., Tsvetkov, Y., Hajishirzi, H., Koh, P.W.,
Zettlemoyer, L.: Spurious Rewards: Rethinking Training Signals in RLVR.
arXiv:2506.10947 (2025)

\bibitem{ref_collapse}
Shumailov, I., Shumaylov, Z., Zhao, Y., Papernot, N., Anderson, R., Gal,
Y.: AI models collapse when trained on recursively generated data. Nature
(2024). \doi{10.1038/s41586-024-07566-y}

\bibitem{ref_restem}
Singh, A., Co-Reyes, J.D., Agarwal, R., et al.: Beyond Human Data: Scaling
Self-Training for Problem-Solving with Language Models. arXiv:2312.06585
(2023)

\bibitem{ref_rest}
Gulcehre, C., Le Paine, T., Srinivasan, S., et al.: Reinforced
Self-Training (ReST) for Language Modeling. arXiv:2308.08998 (2023)

\bibitem{ref_teacherfree}
Strozzi, I.L.: Teacher-Free Self-Training Amplifies but Does Not Compound:
A Pass@k Crossover on a Free-Verifier Domain. arXiv:2606.07856 (2026)

\bibitem{ref_oneshot}
Wang, Y., Yang, Q., Zeng, Z., Ren, L., Liu, L., Peng, B., Cheng, H., He, X.,
Wang, K., Gao, J., Chen, W., Wang, S., Du, S.S., Shen, Y.: Reinforcement
Learning for Reasoning in Large Language Models with One Training Example.
In: NeurIPS (2025). arXiv:2504.20571

\bibitem{ref_phantom}
Xu, C., Yan, N., Chen, L., Kechadi, M.-T.: Phantom Gains: Auditing
Self-Improvement Against a Measured Null. arXiv:2608.20290 (2026)

\bibitem{ref_tinyv}
Xu, Z., Li, Y., Jiang, F., Ramasubramanian, B., Niu, L., Lin, B.,
Poovendran, R.: TinyV: Reducing False Negatives in Verification Improves RL
for LLM Reasoning. arXiv:2505.14625 (2025)

\bibitem{ref_rft}
Yuan, Z., Yuan, H., Li, C., Dong, G., Lu, K., Tan, C., Zhou, C., Zhou, J.:
Scaling Relationship on Learning Mathematical Reasoning with Large Language
Models. arXiv:2308.01825 (2023)

\bibitem{ref_rllimit}
Yue, Y., Chen, Z., Lu, R., Zhao, A., Wang, Z., Yue, Y., Song, S., Huang,
G.: Does Reinforcement Learning Really Incentivize Reasoning Capacity in
LLMs Beyond the Base Model? arXiv:2504.13837 (2025)

\bibitem{ref_star}
Zelikman, E., Wu, Y., Mu, J., Goodman, N.D.: STaR: Bootstrapping Reasoning
With Reasoning. arXiv:2203.14465 (2022)

\bibitem{ref_simplerl}
Zeng, W., Huang, Y., Liu, Q., Liu, W., He, K., Ma, Z., He, J.: SimpleRL-Zoo:
Investigating and Taming Zero Reinforcement Learning for Open Base Models in
the Wild. arXiv:2503.18892 (2025)

\end{thebibliography}

\end{document}
